\subsection{补算：五矩中的$5/3$到底从哪一步出现}
令$\mathsf P=p\mathsf I$，则$U=3p/2$。内能方程的时间和对流部分是
\begin{equation}
 \partial_tU+\grad\cdot(U\vel)
 ={3\over2}Dp+{3\over2}p\theta.
\end{equation}
压力功再提供$p\theta$，所以
$(3/2)Dp+(5/2)p\theta=0$，乘$2/3$才得到
$Dp+(5/3)p\theta=0$。
利用$D\ln n=-\theta$：
\begin{equation}
 D\ln p={5\over3}D\ln n,\qquad D(pn^{-5/3})=0.
\end{equation}
在三个速度自由度的绝热压缩中，密度增大8倍，压强增大$8^{5/3}=32$倍，温度增大4倍。等温闭合则只使压强增大8倍；两者并非同一模型的两种写法。

\subsection{十矩的六个压力方程：看一个分量就能读懂其余分量}
在二维$\partial_z=0$中，$P_{xx}$方程仍包含
\begin{align}
 DP_{xx}+P_{xx}\theta
 +2P_{xx}\partial_xu_x+2P_{xy}\partial_yu_x
 +\partial_xQ_{xxx}+\partial_yQ_{xxy}
 =2{q\over m}(P_{xy}B_z-P_{xz}B_y)+C_{xx}.
\end{align}
$P_{xz}$方程则是
\begin{align}
 DP_{xz}+P_{xz}\theta
 &+P_{xx}\partial_xu_z+P_{xy}\partial_yu_z
 +P_{xz}\partial_xu_x+P_{yz}\partial_yu_x\\
 &+\partial_xQ_{xxz}+\partial_yQ_{xyz}\nonumber\\
 &={q\over m}\{P_{yz}B_z-P_{zz}B_y
                  +P_{xx}B_y-P_{xy}B_x\}+C_{xz}.\nonumber
\end{align}
即使空间变化只在两个方向，$u_z$的剪切、面外压力和磁旋转仍参与动力学。十矩保留的六个对称分量是$xx,yy,zz,xy,xz,yz$，不是九个独立量。

\subsection{局地闭合的确切含义：它可以不搬运热量}
把$\mathsf P=p\mathsf I+\mathsf A$，$\operatorname{tr}\mathsf A=0$，代入局地闭合
$\partial_kQ_{ijk}\approx\nu A_{ij}$。压力方程中闭合导致
\begin{equation}
 \left.\partial_t\mathsf A\right|_{\rm cl}=-\nu\mathsf A,\qquad
 \left.\partial_tp\right|_{\rm cl}=0.
\end{equation}
若均匀无流且没有磁旋转，解为$\mathsf P(t)=p_0\mathsf I+\mathsf A(0)e^{-\nu t}$。例如初始对角压力$(3,1,2)$，其均值$p_0=2$保持不变，三个分量趋向$(2,2,2)$。这是本地各向同性化，没有把能量从热点运到冷点。

因此这种近似虽被称为“热流散度闭合”，其效果与真实非局地热输运并不相同。若想描述沿磁场不同位置之间的热交换，需要空间梯度或沿场非局地算子；仅改变$\nu$不能自动补回这部分信息。\citep{Wang2015,Ng2020,Closure2026}

\subsection{CGL投影的中间计算：平行压力为何多出$2A$}
令$A=b_ib_j\partial_ju_i$、$\theta=\grad\cdot\vel$，
$P_{ij}=p_\perp\delta_{ij}+(p_\parallel-p_\perp)b_ib_j$。
有$b_iP_{ik}=p_\parallel b_k$，所以两个应变项的平行投影分别为
\begin{equation}
 b_ib_jP_{ik}\partial_ku_j=p_\parallel A,\qquad
 b_ib_jP_{jk}\partial_ku_i=p_\parallel A.
\end{equation}
又因$\vct b\cdot D\vct b=0$，
$b_ib_jDP_{ij}=Dp_\parallel$，从而得到
$Dp_\parallel+p_\parallel\theta+2p_\parallel A=0$。

压力方程取迹则给
\begin{equation}
 D(2p_\perp+p_\parallel)
 +(2p_\perp+p_\parallel)\theta
 +2[p_\perp\theta+(p_\parallel-p_\perp)A]=0.
\end{equation}
减去刚才的平行式，除以2，便得到
$Dp_\perp+p_\perp(2\theta-A)=0$。再配合连续性与相应阶次的磁冻结，才可推出双绝热不变量；只使用压力投影本身还不足以得到关于$B$的不变量。

\paragraph{跟算例与适用条件。}
若一个理想冻结、慢变磁通管只作垂直压缩，长度不变，$B\propto n$。CGL给$p_\perp\propto n^2$、$p_\parallel\propto n$，所以$T_\perp\propto n$而$T_\parallel$不变。若一个旋转对称分布被压缩后出现这种各向异性，五矩会强行将其抹平；十矩可以保留，但需要合适闭合。磁零点、非磁化轨道或强热流区域不满足这段推导的前提。
